% SHARED SECTION, \input by both NonGravForces_Theory_and_Implementation.tex and
% NonGravForces_UserGuide.tex. It carries no preamble and opens no document, so that the
% theory exists ONCE and both documents render the same words (R1).
\section{Scope, notation and the force hierarchy}
\subsection{Scope}
\texttt{NonGravForces} computes the non-gravitational acceleration $\vv a_{ng}$, the surface and body torques $\vv\tau$, and their partial derivatives, for a spacecraft described by a panel (macro) model, given its state, attitude and environment. It computes nothing else: time scales, Earth orientation, frame transformations, geodetic coordinates, ephemerides and attitude kinematics are the \texttt{referenceSystem} toolbox; Level-1B reading is \texttt{gracefoReading}; orbit integration is \texttt{gjToolbox}, which \texttt{app.DynamicsSource} serves through that toolbox's frozen dynamics interface. Parameter estimation is NOT delegated: the calibration problem is linear and separates by axis, so \texttt{app.calibrate} carries its own equilibrated QR (owner ruling of 2026-09-09; \texttt{lsqvce} was a dependency until then). This partition (HANDOVER Part~A.2) is the single most important design decision and it is what makes the toolbox smaller than either predecessor while covering more physics.

\subsection{Notation}
$\vv r,\vv v$ position and velocity in GCRS [\si{\metre}, \si{\metre\per\second}]; $\vv q$ attitude quaternion BODY$\to$GCRS (scalar-first Hamilton, passive convention of \texttt{referenceSystem}); $R_{b}$ the rotation matrix; $m$ spacecraft mass; panel $k$ with area $A_k$, outward unit normal $\hat{\vv n}_k$, centroid $\vv c_k$ (body frame), optical triplets $(\alpha,\rho_s,\rho_d)$ in the visible band and $(\varepsilon,\rho_s,\rho_d)$ in the infrared, energy accommodation $\alpha_E$. $\hat{\vv s}$ unit vector to the Sun, $\hat{\vv u}=\vv v_r/|\vv v_r|$ the relative-flow direction, $\theta_k$ the incidence angle with $\gamma_k=\cos\theta_k=\hat{\vv n}_k\!\cdot\!(-\hat{\vv u})$ for drag and $\hat{\vv n}_k\!\cdot\!\hat{\vv s}$ for radiation. SI throughout; angles in radians; $K$ epochs are stored as $3\times K$ arrays.

\subsection{Force hierarchy at the GRACE-FO altitude ($\sim$\SI{490}{\kilo\metre})}
\begin{table}[h]\centering\footnotesize
\begin{tabular}{@{}lllp{4.4cm}@{}}\toprule
Effect & along-track & cross-track & source / note\\\midrule
atmospheric drag & \num{5e-8} (solar min) \ldots\ \num{5e-7} (max) & $\sim\num{e-9}$ & Vielberg \& Kusche 2020; Wöske et al.\ 2019\\
solar radiation pressure & \SI{13}{nm/s^2} RMS & \SI{21}{nm/s^2} & Siemes et al.\ 2023, GRACE-A 2008\\
thermal re-radiation (TRP) & \SI{0.8}{nm/s^2} & \SI{6.4}{nm/s^2} & second-largest radiation term\\
Earth albedo & \SI{0.2}{nm/s^2} & \SI{0.2}{nm/s^2} & $<\SI{1}{nm/s^2}$ combined with IR\\
Earth infrared & \SI{0.2}{nm/s^2} & \SI{0.1}{nm/s^2} & \\
thruster impulses & up to \num{e-6} & & THR1B; Bandikova et al.\ 2019\\
CoM-offset coupling & $\dot{\vv\omega}\times\Delta\vv r+\vv\omega\times(\vv\omega\times\Delta\vv r)$ & & $|\Delta\vv r|\le\SI{100}{\micro\metre}$ (JPL D-56935)\\
antenna thrust, Lorentz & $\lesssim\num{e-11}$ & & hooks only\\
accelerometer noise floor & \SI{1e-10}{m\,s^{-2}\,Hz^{-1/2}} & $\sim10\times$ worse & Kornfeld et al.\ 2019\\\bottomrule
\end{tabular}
\caption{Magnitudes in \si{\metre\per\second\squared} unless stated. Every unit of the toolbox is validated against these orders of magnitude before any finer oracle.}
\end{table}

\section{Atmospheric drag and lift}
\subsection{Free-molecular flow and the per-panel force}
At \SI{400}{}--\SI{500}{\kilo\metre} the Knudsen number is $\gg1$; molecules do not interact with each other near the surface and the force is the integral over the incident Maxwellian of the momentum exchanged at each panel. With mass density $\rho$, relative velocity $\vv v_r$ and per-panel pressure and shear coefficients $C_{p,k}, C_{\tau,k}$,
\begin{equation}
\vv F_{drag}=-\tfrac12\rho\,|\vv v_r|^2\sum_k A_k\left[C_{p,k}\,\hat{\vv n}_k+C_{\tau,k}\,\hat{\vv t}_k\right],\qquad
\hat{\vv t}_k=\frac{\hat{\vv u}-(\hat{\vv u}\!\cdot\!\hat{\vv n}_k)\hat{\vv n}_k}{|\hat{\vv u}-(\hat{\vv u}\!\cdot\!\hat{\vv n}_k)\hat{\vv n}_k|},
\label{eq:drag}
\end{equation}
where the sign convention makes $C_p>0$ push the panel inward along $-\hat{\vv n}_k$ for a windward panel. The projection onto $\hat{\vv u}$ is drag, the remainder is lift plus side force. The kernel returns the per-panel $\vv F_k$ (principle P2); $\vv F$ and $\vv\tau=\sum_k(\vv c_k-\vv r_{cm})\times\vv F_k$ are reductions.

\subsection{Relative velocity}
\begin{equation}
\vv v_r=\vv v-\vv\omega_\oplus\times\vv r-\vv v_{wind},
\end{equation}
with the co-rotation term formed in \texttt{referenceSystem.applyRotation} (it carries $\vv\omega\times\vv r$ exactly, never by hand) and $\vv v_{wind}$ from HWM14 (Drob et al.\ 2015; quiet-time RMS \SI{40}{}--\SI{80}{\metre\per\second}) as an opt-in. At \SI{7.6}{\kilo\metre\per\second} the co-rotation term is \SI{\sim470}{\metre\per\second} at the equator and enters cross-track drag at the \SI{6}{\percent} level; winds of \SI{100}{\metre\per\second} enter at \SI{1.3}{\percent} of the along-track force but dominate the cross-track signal (Doornbos 2011).

\subsection{Sentman's model (exact form implemented)}
Let $s=|\vv v_r|/\sqrt{2\mathcal RT_\infty/M}$ be the speed ratio with $M$ the mean molar mass of the local gas (\emph{not} air: at \SI{500}{\kilo\metre} atomic oxygen with $M\approx\SI{16}{g/mol}$ dominates, so $s\approx8$; using $R_{air}=\SI{287}{J\,kg^{-1}\,K^{-1}}$ as v1 did biases $s$ by $\sqrt{29/16}=1.35$), $\gamma=\cos\theta$, $\ell=\sin\theta$, and
\begin{equation}
P=\frac{e^{-\gamma^2s^2}}{s},\quad Q=1+\frac1{2s^2},\quad G=\frac1{2s^2},\quad Z=1+\operatorname{erf}(\gamma s),\quad
\frac{v_{re}}{v_{in}}=\sqrt{\tfrac12\Big[1+\alpha_E\Big(\frac{4\mathcal RT_w}{M|\vv v_r|^2}-1\Big)\Big]}.
\end{equation}
\textbf{Which pair is which.} Sentman's published coefficients, and Doornbos (2011, eqs 3.55--3.56), are \emph{drag} and \emph{lift} --- along and across $\vv v_r$:
\begin{align}
C_D&=\frac{P}{\sqrt\pi}+\gamma QZ+\frac{\gamma}{2}\frac{v_{re}}{v_{in}}\big(\gamma\sqrt\pi Z+P\big),\label{eq:sentmanCD}\\
C_L&=\ell GZ+\frac{\ell}{2}\frac{v_{re}}{v_{in}}\big(\gamma\sqrt\pi Z+P\big).\label{eq:sentmanCL}
\end{align}
The toolbox does not use that pair. Its kernel is written in the plate's own frame, \emph{normal} and \emph{tangential} (eq.~\eqref{eq:drag}), which is what a per-panel force and a torque both need, and it obtains those from the same momentum integrals:
\begin{align}
C_p&=\underbrace{\frac{\gamma P}{\sqrt\pi}+\Big(\gamma^2+\frac1{2s^2}\Big)Z}_{\text{incident}}+\underbrace{\frac12\frac{v_{re}}{v_{in}}\big(\gamma\sqrt\pi Z+P\big)}_{\text{re-emitted}},\label{eq:sentmanCp}\\
C_\tau&=\ell\Big(\frac{P}{\sqrt\pi}+\gamma Z\Big).\label{eq:sentmanCtau}
\end{align}
The two pairs are one rotation apart, $C_D=C_p\gamma+C_\tau\ell$ and $C_L=C_p\ell-C_\tau\gamma$, and the toolbox asserts both projections against eqs.~\eqref{eq:sentmanCD}--\eqref{eq:sentmanCL}. \textbf{This document labelled eqs.~\eqref{eq:sentmanCD}--\eqref{eq:sentmanCL} $C_p$ and $C_\tau$ until 0.10.0}, which is the same confusion that produced defect D1 below; nothing in the code ever read it, and it is corrected here. Note in particular that the \emph{tangential} coefficient carries no $1/s^2$ term: the lone $Z/(2s^2)$ belongs to the lift. The limit that settles it is edge-on ($\gamma\to0$): $C_p\to1/(2s^2)$ and $C_\tau\to\ell/(s\sqrt\pi)$, and the hyperthermal ram limit is $C_p\to2\gamma^2$, $C_\tau\to2\gamma\ell$, $C_D\to2\gamma$.

\paragraph{How good is the closed form?} Incomplete accommodation is treated in three ways in the
literature, and this toolbox implements the first: (a) the energy-flux form above, where
$v_{re}/v_{in}$ carries $(1-\alpha_E)/2$ of the incident energy; (b) the kinetic-temperature form
$T_{k,r}=T_{k,i}(1-\alpha_E)+\alpha_E T_w$ with $T_{k,i}=mU^2/3k$, which carries $(1-\alpha_E)\,2/3$
(Walker et al.\ 2014, eqs 10--12); and (c) sampling the Cercignani--Lampis--Lord kernels, which is the
physically correct reflected velocity distribution and requires a DSMC rather than a closed form.
At $\alpha_E=1$ (a) and (b) are the SAME NUMBER, and eq.~\eqref{eq:sentmanCp} reproduces Walker's
published flat-plate coefficient to \num{1.6e-15} there --- the only check of this chain against
arithmetic written outside this project. Below full accommodation they separate: measured on the
plate, (a) reads \SI{1.7}{\percent} below (b) at $\alpha_E=0.95$, \SI{2.4}{\percent} at $0.90$ and
\SI{5.6}{\percent} at $\alpha_E=0$. Walker et al.\ measure (b) against (c) for the GRACE geometry at
\SI{\sim1.4}{\percent} for $\alpha_E=0.9$ and \SI{\sim1.7}{\percent} at $\alpha_E=0$. Both are of order
\SIrange{1}{2}{\percent} of $C_D$ at the $\alpha_E\approx0.85$--$1$ of a real LEO surface, inside the
\SIrange{5}{10}{\percent} that $\alpha_E$ itself is uncertain by --- so the closed form is kept, and
the deviation is stated instead of assumed to be zero.

Diffuse re-emission is isotropic about the normal, so it contributes to $C_p$ and \emph{not} to $C_\tau$ --- which is why $C_\tau$ carries no $v_{re}/v_{in}$ while $C_L$ does. The re-emission ratio carries the wall temperature $T_w$ and the accommodation coefficient $\alpha_E$; complete accommodation ($\alpha_E=1$) gives diffuse re-emission at $T_w$. \textbf{The gas force is therefore not antiparallel to $\vv v_r$}: $C_L\ne0$ off the ram direction, \SI{0.34}{\degree} on GRACE-FO (measured). Only \texttt{constantCd} is antiparallel, by construction. \textbf{Defect D1 of v1}: the exponential group and the whole shear term were divided by $s$ instead of $s^2$, an error of $s\approx8$ in the shear on inclined panels; the toolbox pins this as a regression test with the correct value from an independent kinetic-theory integration (\texttt{tools/pyxcheck/sentman\_check.py}).

\subsection{Other gas--surface interaction models}
\emph{Maxwell}: a fraction $\sigma$ diffuse, $1-\sigma$ specular. \emph{Cercignani--Lampis--Lord (CLL)}: separate normal and tangential accommodation $(\alpha_n,\sigma_t)$ with quasi-specular lobes; used by Walker et al.\ (2014) and Mehta et al.\ (2014). \emph{Schamberg}: quasi-specular lobe of adjustable width. All are \texttt{drag.gsi} options; none is in a preset except Sentman, because the density error (\S\ref{sec:density}) exceeds their differences except below \SI{300}{\kilo\metre}.

\subsection{Energy accommodation}
Published values: March et al.\ (2021) optimise a \emph{constant} $\alpha_E$ per mission against attitude manoeuvres and wind analyses and obtain $0.85$ (CHAMP) and $0.82$ (GOCE), recommending $\alpha_E=0.85$ and superseding the $0.93$ of Doornbos (2011); $0.89$ from CHAMP yaw (Pilinski \& Argrow 2013). They fit rather than compute \emph{because} the in-situ composition and temperature a physical model needs are not measured on any of these satellites. The SESAM model (Pilinski et al.\ 2010, 2013) makes $\alpha_E$ a function of the atomic-oxygen partial pressure $p_O$ through a Langmuir isotherm $\Theta=Kp_O/(1+Kp_O)$, $\alpha_E=\Theta\,\alpha_{ads}+(1-\Theta)\,\alpha_{clean}$, reproducing measurements to \SI{\sim3}{\percent}. The toolbox implements \ngfoptdragaccommodation. \texttt{pilinski2010} is the Langmuir isotherm in $n_O T$ (Pilinski 2011, \S7.3); \texttt{walker2014} carries the adsorbate coverage with a clean-substrate $\alpha_s=3\mu/(1+\mu)^2$ (Paul et al.\ 2022, eqs 10--13). \texttt{sesam} is the model of Pilinski et al.\ (2013) itself: the same mixture, but with $K_L$ a function of the incident kinetic energy through an initial sticking coefficient (their eq.~26b, the fraction of the incident Maxwellian falling below a binding energy of \SI{5.7}{\electronvolt}), and with $P_O$ a \emph{ram} pressure $\tfrac12\rho_O V_r^2 C_{D,\mathrm{in},O}$ rather than $n_O k T$ --- at \SI{520}{\kilo\metre} a factor of 140, so the two parameterisations' Langmuir constants are not interchangeable either. Its substrate coefficient uses Goodman's $2.4\mu/(1+\mu)^2$ where \texttt{walker2014} ships $3.0$; the sources disagree and each option follows its own. \textbf{Not implemented from that paper}: the DSMC lookup tables and the tumbling-sphere shape factor, which map $\alpha$ to a drag coefficient --- this toolbox has its own panel geometry and Sentman kernel and takes only $\alpha$ from SESAM (HANDOVER O-16). A document that offers an option the code refuses is worse than one that offers none, which is why this list is generated from \texttt{Options.validateValue}.

\textbf{All three presets take the fitted $\alpha_E=0.85$, and the adsorption models are opt-ins.} That is a ruling (HANDOVER O-42, 2026-09-17), not an approximation, and the reason is measured. Evaluated at CHAMP's altitude and CHAMP's solar activity (\SI{300}{}--\SI{350}{\kilo\metre}, $F_{10.7}=70$), \texttt{pilinski2010} returns $0.86$--$0.95$, straddling the $0.85$--$0.93$ March et al.\ (2021) report for SESAM there: the adsorption family reproduces the literature where the literature was fitted. The same code returns $0.11$--$0.17$ at GRACE-FO's \SI{520}{\kilo\metre} in deep solar minimum, and $0.63$--$0.83$ at that altitude at $F_{10.7}=180$, because the atomic-oxygen partial pressure driving $\Theta$ falls by two orders of magnitude between them ($\Theta=0.83$ at \SI{300}{\kilo\metre}, $0.05$ at \SI{520}{\kilo\metre}). So these models predict a factor-of-five solar-cycle swing in GRACE-FO's accommodation, and \SI{2020}{}-01-01---the day this toolbox is validated on---sits at its extreme low end. On that day the choice is worth \SI{17.5}{\percent} of the drag (\texttt{walker2014} \SI{8.09e-9}{\metre\per\second\squared} against \SI{6.88e-9}{} for $\alpha_E=0.85$; $0.93$ gives \SI{6.62e-9}{}), which for accelerometer calibration is absorbed by the estimated scale factor and propagates into any inverted density. Both positions are defensible; the presets take the one the missions are actually processed with, and the other is one option away.

\subsection{Density and composition}\label{sec:density}
Backends, as \texttt{drag.density} admits them: \ngfoptdragdensity\ --- exponential (sizing only), NRLMSISE-00 (Picone et al.\ 2002; the Brodowski C port, whose 17-case DOCUMENTATION table is the bit-exact oracle) and NRLMSIS~2.1 (Emmert et al.\ 2021, NRL's unmodified Fortran, checked against its own 200 reference cases). DTM2020 (Bruinsma \& Boniface 2021) and JB2008 (Bowman et al.\ 2008) are \emph{not} implemented and were removed from the option set in 0.10.0, because an option that cannot be evaluated is a promise the code does not keep; their space-weather drivers stay readable in \texttt{io.spaceWeather}. Inputs: $F_{10.7}$ (\emph{observed}, previous day, and 81-day mean; v1 fed adjusted same-day values, $\pm\SI{3.4}{\percent}$), $F_{30}$, $S_{10.7}, M_{10.7}, Y_{10.7}$ (JB2008, from SET), $A_p$ 3-hourly history (the \texttt{ap7} vector), $K_p$/Hp, $D_{st}$ for storms. Uncertainty: \SI{10}{}--\SI{30}{\percent} between models and data sets; NRLMSIS~2.1 is \SI{\sim10}{\percent} below NRLMSISE-00 at GRACE altitudes (Guo \& Zhang 2024). Composition is required by the GSI ($M$, and $p_O$ for SESAM). The models are evaluated at the \emph{geodetic} height on WGS-84: a spherical Earth gives \SI{-21.4}{\kilo\metre} at the poles and a density factor of 1.40 on a polar mission (v1 defect D2); the conversion is \texttt{referenceSystem.cartesianToGeodetic} with the ellipsoid passed explicitly (its default is GRS80).

\subsection{Macro model and self-shadowing}
GRACE-FO (JPL D-56935 Table~5): eight plates, side panels at \SI{50}{\degree} with normals printed as $\sin50^\circ=0.766044$, $\cos50^\circ=0.642787$ (their squares sum to $1-\num{5e-7}$; the toolbox accepts $10^{-5}$ and documents it), total area \SI{16.9163284}{\metre\squared}, mass \SI{601.214}{kg} (C) / \SI{601.209}{kg} (D), inertia tensor Table~4. The body is not convex: the flanks expose their undersides, so $\sum_kA_k\hat{\vv n}_k\ne0$ and self-shadowing is physical. Wöske et al.\ (2019) quote $\sim\SI{10}{\percent}$ force error of an 8-plate model against finite-element ray tracing. Self-shadowing options (\texttt{geom.selfShadow}): \ngfoptgeomselfShadow\ --- none, an equiangular CUBE-MAP look-up over incidence directions with $O(1)$ bilinear retrieval (\texttt{geom.lutResolution} cells per cube-face edge, convergence-tested), or an exact ray cast. The table records OCCLUSION ONLY: whether a plate faces the source is the kernels' business, and folding the two together is what deleted the edge-on shear before 0.10.0 (\S\ref{sec:flow}). The S-band boom is a cylinder (projection area \SI{0.0461901}{\metre\squared}, \SI{0.27}{\percent}); it is a separate element kind, not a plate.

\subsection{Occlusion is not facing, and a gas is not a beam}\label{sec:flow}
A look-up table that returns $f=0$ for a plate the source cannot see is a rule about a PARALLEL BEAM.
Applied to the ambient gas it is wrong twice, and both errors shipped until 0.10.0 (review findings
F1 and F2). First, a thermal gas arrives from every direction of the front hemisphere with the
drifting-Maxwellian weight, so an edge-on plate ($\gamma=0$) still receives flux and still carries
Sentman's shear $C_\tau\to\ell/(s\sqrt\pi)$ --- setting it to zero deleted that shear on
$\sim\SI{15}{\metre\squared}$ of GRACE-FO plate. Second, the same column was being used for the
Earth's radiation, which zeroed the heating of every panel facing away from the \emph{Sun}, and the
nadir panel's steady-state temperature read \SI{0}{\kelvin} at 11 of 120 real epochs.

The separation is now explicit: \texttt{geom.selfShadow} answers only ``is anything in the way'', and
whether a plate faces a source is the kernels' business. For the gas, \texttt{geom.flowShadow}
computes per plate, epoch and species three ratios of occlusion-weighted to unweighted directional
integrals of the incident flux over the front hemisphere --- number flux (for the re-emitted
momentum), normal momentum and tangential momentum --- with closed-form speed moments and a
flow-aligned product quadrature over directions. The limits that pin it: a plate with nothing in
front of it returns exactly 1, bitwise; and $s\to\infty$ reduces to the parallel-beam occlusion along
$\vv v_r$, which is what \texttt{constantCd} uses. Not modelled, and bounded by the oracle:
re-emitted molecules striking a second plate. Measured on 120 real epochs, the correction moves the
drag by \SI{1.68e-9}{\metre\per\second\squared} RMS, almost all along track, and the thermal recoil
by \SI{1.77e-9}{\metre\per\second\squared} RMS, radial.

\subsection{Implementation notes}
Vectorise over $K$ epochs: $\gamma$ is a $N\times K$ array, $\operatorname{erf}$ and $\exp$ are evaluated once per array; the density backend accepts $K$ points and returns $\rho$, species number densities and $T_\infty$; the MSIS inner loop is the one place where hand-written C (MEX) is sanctioned for post-processing. NaN density is an error (\texttt{numericNaN}), never a zero force.

\section{Solar radiation pressure}
\subsection{Irradiance and pressure}
Total solar irradiance $S_0=\SI{1361}{W/m^2}$ at \SI{1}{au} (IAU 2015 Resolution B3; SORCE/TIM; TSIS-1: $1361.5\pm5.9$), $P_0=S_0/c=\SI{4.5398e-6}{Pa}$. The older \SI{1367}{W/m^2} (Behzadpour et al.\ 2021, Table~1) is available as \texttt{Options.srp.tsi} for legacy comparisons and is never a constant (v1 used both values in two routines, \SI{0.44}{\percent} inconsistency; v0.1.1 repeated the class of error between its oracle and its constant). Scale with $(\mathrm{au}/|\vv r_s-\vv r|)^2$; the Sun position comes from \texttt{referenceSystem.geocentricBodyState} (analytic SOFA tier, RMS \SI{8.3e3}{\metre}; DE440 tier through MICE), never from a second implementation.

\subsection{Per-panel optical model}
For a panel with $\gamma=\hat{\vv n}_k\!\cdot\!\hat{\vv s}>0$,
\begin{equation}
\vv F_k=-\nu\,P\,A_k\,\gamma\Big[(\alpha+\rho_d)\,\hat{\vv s}+\Big(2\rho_s\gamma+\tfrac23\rho_d\Big)\hat{\vv n}_k\Big],\qquad \alpha+\rho_s+\rho_d=1,
\label{eq:srp}
\end{equation}
with $\nu\in[0,1]$ the shadow function and $2/3$ the Lambertian recoil factor. The same kernel serves Earth albedo (visible triplet) and Earth infrared (infrared triplet: $\varepsilon,\rho_s,\rho_d$) through the \emph{source list} $\{\hat{\vv d}_j,P_j,\text{band}_j\}$ (principle P3): SRP is one source, point-source Earth radiation two, a tiled Earth $2M$. Thermal re-emission of the absorbed part is \emph{not} folded into eq.~\eqref{eq:srp} by default (\S\ref{sec:trp}).

\subsection{Shadow function}
Apparent radii $a=\arcsin(R_s/|\vv r_s-\vv r|)$, $b=\arcsin(R_\oplus/|\vv r|)$, angular separation $c$ between $-\hat{\vv r}$ and $\widehat{\vv r_s-\vv r}$. Cylindrical model: $\nu=0$ iff $\vv r\!\cdot\!\hat{\vv s}<0$ and $|\vv r-(\vv r\!\cdot\!\hat{\vv s})\hat{\vv s}|<R_\oplus$. Conical model (Montenbruck \& Gill 2000, \S3.4.2): full light for $c\ge a+b$, umbra for $c\le b-a$, annular ($\nu=1-b^2/a^2$) for $c\le a-b$, else
\begin{equation}
x=\frac{c^2+a^2-b^2}{2c},\quad y=\sqrt{a^2-x^2},\quad
A=a^2\arccos\frac xa+b^2\arccos\frac{c-x}{b}-cy,\quad \nu=1-\frac{A}{\pi a^2}.
\end{equation}
The conical function is continuous but only $C^0$ at the umbra edge; for orbit integration a $C^1$ transition is required or the integrator collapses its step. SOLAARS-CF (Robertson 2015; Zhang et al.\ 2019) adds Earth oblateness and atmospheric extinction/refraction as a tabulated attenuation curve; Siemes et al.\ (2023) quantify the model-version effect at \SI{6}{}--\SI{9}{\percent} of the radiation-pressure acceleration. Lunar eclipses use the same geometry with $R_{moon}$ and are an opt-in. Tests: $\nu\in[0,1]$, monotone across the penumbra, symmetric in the geometry, and the worked example of Montenbruck \& Gill; a numerical disc integration in Python is the independent oracle.

\section{Earth radiation pressure}
\subsection{Knocke--Ries--Tapley (1988)}
Albedo and emissivity as zonal series with an annual term, $t_0$ = 22 December:
\begin{equation}
a(\phi,t)=a_0+a_1\cos\omega(t-t_0)\,P_1(\sin\phi)+a_2P_2(\sin\phi),\quad
e(\phi,t)=e_0+e_1\cos\omega(t-t_0)\,P_1(\sin\phi)+e_2P_2(\sin\phi),
\end{equation}
$a_0=0.34,a_1=0.10,a_2=0.29$; $e_0=0.68,e_1=-0.07,e_2=-0.18$; $\omega=2\pi/\SI{365.25}{d}$. Each visible surface element $dA$ at distance $\varrho$ with outgoing angle $\vartheta$ is a Lambertian source: albedo flux $dP=\dfrac{a\,S_0\cos\theta_s\,\cos\vartheta\,dA}{\pi c\,\varrho^2}$ (only for $\cos\theta_s>0$, gated by the shadow function of the element), infrared flux $dP=\dfrac{e\,S_0\cos\vartheta\,dA}{4\pi c\,\varrho^2}$. The visible cap is tiled ($12\times24$ in \texttt{standard}, $24\times48$ in \texttt{precise}); a few hundred tiles suffice at LEO (convergence test).
\subsection{CERES}
Observed top-of-atmosphere fluxes (EBAF climatology, SYN1deg hourly) replace the zonal model at the nm/s$^2$ level (Vielberg \& Kusche 2020; Wang et al.\ 2023); anisotropic reflection through the CERES angular distribution models (Bhattarai et al.\ 2019) is \emph{not} modelled: Lambertian reflection is the model, stated in the unit's ACCURACY block. An \texttt{erp.anisotropy} option was accepted and silently ignored until 0.10.0 and has been removed (review finding F4). The backends are \ngfopterpbackend. The netCDF variable names must be read from a real file (open item O-3).

\section{Thermal re-radiation}\label{sec:trp}
A panel at temperature $T_k$ with emissivity $\varepsilon_k$ emits $\varepsilon_k\sigma T_k^4A_k$ and recoils with
$\vv F_k=-\tfrac23\,\varepsilon_k\sigma T_k^4A_k\,\hat{\vv n}_k/c$. In steady state $T_k$ follows from the balance $\alpha_kQ_{sun,k}+\alpha_kQ_{alb,k}+\varepsilon_kQ_{ir,k}+q_{int,k}=\varepsilon_k\sigma T_k^4$; the transient model adds $C_k\dot T_k$ with a per-panel heat capacity (Wöske 2019 thesis, Ch.~4; Rievers \& Lämmerzahl 2011). For GRACE, TRP is the second-largest radiation term after SRP (Wöske et al.\ 2019; Vielberg \& Kusche 2020; Siemes et al.\ 2023: \SI{6.4}{nm/s^2} cross-track RMS), so ``folding it into SRP'' (the $2/3$ factor on absorbed flux) is a documented, deliberately crude option, not the default. Metamorphic test: in steady state the total absorbed power equals the total emitted power.

\section{Other effects and artefacts}
\textbf{Thruster forces}: cold-gas attitude thrusters (THR1B) produce linear accelerations up to \SI{e-6}{\metre\per\second\squared} that survive the \SI{35}{mHz} CRN filter (Meyer et al.\ 2012; Bandikova et al.\ 2019 model them as impulse responses per thruster pair). The toolbox models them as impulses with the CRN response and uses THR1B epochs as \emph{priors} for the impulse detector (\S\ref{sec:qc}). \textbf{Antenna thrust} (K-band, GPS): $\lesssim\SI{e-11}{\metre\per\second\squared}$, hook only. \textbf{Lorentz force}: negligible at GRACE-FO, hook only. \textbf{Artefacts the model does not explain}: twangs (up to \SI{2e-5}{\metre\per\second\squared} radial; Flury et al.\ 2008; Peterseim 2014), heater switching, magnetorquer spikes (Peterseim et al.\ 2012). These are removed by quality control, never modelled.

\section{Torques and the centre-of-mass coupling}
Surface torque $\vv\tau=\sum_k(\vv c_k-\vv r_{cm})\times\vv F_k$ from the same per-panel array (v1's separate torque implementation lost the vis/IR split: \SI{8}{\percent}). Gravity gradient $\vv\tau_{gg}=\frac{3\mu}{r^3}\hat{\vv r}\times(\mathbf I\hat{\vv r})$ with $\hat{\vv r}$ in the body frame. Magnetic $\vv\tau_{m}=\vv m\times\vv B$ with residual and commanded dipoles (MAG1B) and $\vv B$ from IGRF-14 (IAGA 2024; degree 13 main field, degree 8 secular variation; oracle: the igrf14 test cases) evaluated in ITRS and rotated to the body frame.

The accelerometer proof mass sits at $\Delta\vv r$ from the centre of mass ($\le\SI{100}{\micro\metre}$ per axis by design, \SI{\sim50}{\micro\metre} achieved; GRACE-C CoM calibration agrees with JPL to \SI{0.5}{\micro\metre}, Hu et al.\ 2022) and senses
\begin{equation}
\vv a_{sensed}=\vv a_{ng}+\dot{\vv\omega}\times\Delta\vv r+\vv\omega\times(\vv\omega\times\Delta\vv r),
\end{equation}
with $\vv\omega,\dot{\vv\omega}$ from SCA1B through \texttt{referenceSystem.angularVelocityFromRotationRate} after smoothing (Klinger \& Mayer-Gürr 2016; Wöske 2019). This is a property of the \emph{instrument}, not of the physics, and lives in the sensor layer (\S\ref{sec:sensors}).

\section{GRACE-FO data and the measurement model}\label{sec:sensors}
\subsection{Products and frames}
ACC1A/ACT1A (\SI{10}{Hz}) $\to$ ACC1B/ACT1B (\SI{1}{Hz}) via time-tag correction and the CRN low-pass filter (\SI{\sim35}{mHz}, Wu et al.\ 2006); ACT (McCullough et al.\ 2019) substitutes thruster spikes and, for GRACE-D from 2018-06-21, transplants GRACE-C data (Bandikova et al.\ 2019); ACH (2022) hybridises. SCA1B quaternions are ICRF$\to$SRF; the toolbox stores BODY$\to$GCRS and conjugates exactly once in the attitude provider, verified by physics (nadir and line-of-sight), never by a flag. GNV1B is Earth-fixed (ITRS) and is rotated with \texttt{rotationItrsToGcrs} (v1 had no such path and treated GPS tags as UT1: \SI{18}{\second} of Earth rotation). MAS1B holds cold-gas tank masses (\SI{\sim30}{kg}), not spacecraft mass; the handbook constant is the default.

\subsection{Calibration model}
$\vv a_{cal}=\vv b(t)+\mathbf S\,\vv a_{meas}$ with a bias vector (cubic B-spline, 6-hour knots, Klinger \& Mayer-Gürr 2016) and a full $3\times3$ scale/coupling matrix; typical scales $0.83$--$1.05$ (GRACE-C along-track $1.05\pm\num{2e-4}$, Siemes et al.\ 2023), biases $\sim\SI{e-6}{\metre\per\second\squared}$ with thermal excursions of hundreds of nm/s$^2$. Observability: along-track bias/scale well determined (\SI{0.5}{nm/s^2}, 1--3\,\%), cross-track poorly at solar minimum (\SI{22}{nm/s^2}, \SI{\sim15}{\percent}; van Helleputte 2009). The estimation is delegated to \texttt{lsqvce} (ARMA whitening from the sensor noise model, iterative snooping, variance components); \texttt{NonGravForces.app.calibrate} builds the observation vector and the design blocks, including Heaviside columns for detected jumps.

\subsection{The sensor abstraction}
The physics core returns the true dynamics record $\{\vv a_{ng},\vv\tau,\vv\omega,\dot{\vv\omega},t\}$; a \texttt{Sensor} maps it to an observable and provides the partials with respect to its own parameters and its one-sided noise ASD. \textbf{Electrostatic}: bias, scale matrix, CoM coupling, CRN filter; ASD $10^{-10}\sqrt{1+(f/0.5)^4+0.005/f}\ \si{m\,s^{-2}\,Hz^{-1/2}}$ (Kornfeld et al.\ 2019). \textbf{Atom interferometer} (CARIOQA pathfinder, MAGIC/NGGM hybrids): phase $\Delta\phi=k_{eff}\,a\,T^2$ with $k_{eff}=4\pi/\lambda$, interrogation time $T$ (\SI{1}{\second}), cycle time (\SI{4}{\second}, dead time), phase ambiguity modulo $2\pi$ requiring hybridisation with an electrostatic channel (Zahzam et al.\ 2022; Trimeche et al.\ 2019), and a rotation input because nadir pointing implies mrad/s rates whose Coriolis phase must be compensated (Lévèque et al.\ 2022); projection-noise floor \SI{e-10}{m\,s^{-2}\,Hz^{-1/2}}, target \num{e-12}. \textbf{Drag-free}: residual acceleration after control (LISA-Pathfinder heritage; NGGM \SI{e-11}{m\,s^{-2}\,Hz^{-1/2}} class). Optomechanical or MEMS instruments are further subclasses; calibration and quality control are written against the abstract interface, so no core code changes.

\section{Partial derivatives}
For orbit determination and calibration the toolbox returns $\partial\vv a/\partial\vv r$ (through $\partial\rho/\partial h$ and the shadow function), $\partial\vv a/\partial\vv v$ (drag, exact from eq.~\eqref{eq:drag} since $C_p,C_\tau$ depend on $\vv v_r$ through $s$ and $\gamma$), and parameter partials $\partial\vv a/\partial C_D=\vv a_{drag}/C_D$, $\partial\vv a/\partial C_R$, $\partial\vv a/\partial\alpha_E$. Analytic where clean; \texttt{partials.method} admits \ngfoptpartialsmethod\ and nothing else. \textbf{Complex-step differentiation is impossible in this chain and was removed}: MATLAB's \texttt{erf} refuses a complex argument and \texttt{abs}/\texttt{vecnorm} discard the imaginary part, so no complex step survives Sentman's coefficients or any norm on the way. The verifier is Richardson-extrapolated central differences instead, and the accuracy claim is stated against that. One evaluation serves both consumers (\texttt{Model.partials}).

\section{Quality control: outliers, impulses, jumps}\label{sec:qc}
Three event classes with different signatures need three detectors, all operating on whitened residuals $r(t)=y_{meas}-(b+\mathbf S\,y_{model})$ from a preliminary calibration:
\begin{itemize}[nosep]
\item \textbf{Outliers} (isolated samples): Hampel filter, moving median and $1.4826\cdot$MAD over a \SI{61}{\second} window, threshold $4\sigma$; robust to 50\,\% contamination, so thruster spikes and twangs cannot inflate the scale (a moving mean/std would). Implemented and verified against an independent numpy oracle (\texttt{qc.hampel}).
\item \textbf{Impulses} (CRN-shaped pulses, \SI{\sim30}{\second} wide, simultaneous in three axes): matched filter with the CRN impulse response; THR1B epochs are assigned first, unassigned pulses are twang/impact candidates; an impact (not a twang) also changes $\vv\omega$ in SCA1B.
\item \textbf{Jumps} (steps in level or slope): PELT change-point detection (Killick et al.\ 2012) with BIC penalty on piecewise-constant and piecewise-linear models; eclipse entries/exits (from the shadow function) and mass-trim events (MAS1B) are priors. Jumps are \emph{estimated} as Heaviside columns in the calibration, never cut.
\end{itemize}
Output: a canonical event table (\texttt{qc.eventTable}: $t_{start},t_{end}$, class, axis, amplitude, $\sigma$, source, confidence) and a logical mask. Limits: impulses shorter than the CRN width are unresolvable in L1B (use ACC1A); jumps below \SI{\sim1}{nm/s^2} are indistinguishable from $1/f$ bias drift without temperature data; along-track thresholds must adapt to geomagnetic activity because the residual carries the 10--30\,\% density error. Metamorphic tests: time shift, affine scaling, time reversal, planted events recovered with amplitude to $3\sigma$.

\section{Implementation}
\subsection{Architecture}
Nine layers (L0 core, L1 io, L2 env, L3 geom, L4 kernel, L5 \texttt{Model}, L6 sensor, L7 qc, L8 app), \ngfunitcount{} units, one catalogue (\texttt{doc/functionlist.csv}) that generates the function reference and drives the stub and audit tests. A layer calls only lower layers plus the four dependencies. \texttt{Model} is the single point of access (\texttt{accel | torque | series | partials | budget}; \texttt{rhs} was removed in LP-C---an integrator binds \texttt{app.DynamicsSource} instead); \texttt{Options} is a value class with presets \texttt{fast/standard/precise} and validated dotted overrides; extended models are opt-ins that no preset enables, so preset results never move when a model lands.

\subsection{Workflow}
\begin{lstlisting}
run('.../referenceSystem/setupReferenceSystem.m'); run('.../gracefoReading/setupGraceFOReading.m')
setupNonGravForces; runNgfTests
[ep,r,v]  = NonGravForces.io.readGNV1B(gnv);      % ITRS -> GCRS via referenceSystem
[~,q]     = NonGravForces.io.readSCA1B(sca);      % ICRF->SRF conjugated once
[~,y]     = NonGravForces.io.readACC1B(act);
macro = NonGravForces.geom.graceFO("C");
opt   = NonGravForces.Options("standard","drag.wind","hwm14");
M     = NonGravForces.Model(macro, opt, 'Attitude', att, 'Mass', mass);
T     = M.series(ep, r, v, q);                    % a_ng, tau, provenance
sens  = NonGravForces.sensor.Electrostatic(...);
[mask, events] = NonGravForces.qc.screen(y, sens.observe(T), thr, eclipse, mass);
R = NonGravForces.app.calibrate(y, T, sens, 'Events', events, 'Mask', mask);
\end{lstlisting}

\subsection{Error handling}
Every error passes through \texttt{NonGravForces.fail(unit, mnemonic, message)}; the identifier is \texttt{NonGravForces:<unit>:<mnemonic>}, the mnemonic's first token is one of \texttt{input, dep, data, numeric, unsupported, notImplemented}, and the message ends with a correction hint after ``\texttt{ -> }''; \texttt{fail} refuses anything else. Contracts are \texttt{arguments} blocks (shape, class, range) plus semantic checks at the top of the body; nothing is validated twice (the macro model validates once, kernels trust its packed arrays). No fidelity fallback exists; unknown option fields error; absent optional dependencies error at the point of use naming their setup function; dependency errors propagate unwrapped. Every result carries \texttt{info.provenance} (toolbox and dependency versions, option hash, files, warnings, UTC time). NaN policy is stated per unit: a NaN density is an error because it would silently zero the drag; a NaN sample in quality control is never flagged and never poisons a neighbour. Numerical guards report their diagnostic (condition number, residual, epoch index). Tests assert identifiers, not messages, and pair every refusal with the case that must succeed.

\subsection{Testing}
Seven categories per public unit, counted mechanically: contract, reference (an authority outside the project: published tables, independent Python, real data), precision (a bound with the reason for its size in a comment), vectorization ($K$ at once equals one at a time, and chunked/parallel series are bit-identical), robustness (NaN, gaps, empty, flags), metamorphic (rotation covariance, affine/shift/reversal invariance, mass and density scaling, $\nu\in[0,1]$ monotone), speed (ratios against a built-in, absolute times recorded). A passing bound leaves its number in \texttt{tests/last\_run.txt}; the gate is failures + incomplete + audit shortfall + written log, not the pass count. This session: 74 tests, GREEN in R2026a Update 4; two failures were authoring errors found by decomposition (a typed oracle value; a 0/1-based index), one was a guessed bound corrected by measurement (interpreted loop 382$\times$ slower than \texttt{movmedian}; algorithmically equivalent).

\subsection{Performance}
$3\times K$ column layout; environment evaluated in bulk; \texttt{pagemtimes} for attitude; the panel loop over $N\approx8$ with $K$-vectorised arithmetic; self-shadow LUT instead of per-epoch ray casting; \texttt{series} chunks $K$ and optionally runs the chunks under \texttt{parfor} (Parallel Computing Toolbox, feature-detected) with bit-identical results either way; the single-epoch C/MEX kernel exists only for the integrator right-hand side, where MATLAB call overhead (5--15\,\si{\micro\second}) dominates, and links \texttt{referenceSystem}'s \texttt{rsApiRotation}; no GPU (eight panels give a GPU nothing to do).

\section{Audit of the predecessors (summary; full table in HANDOVER Part C)}
\textbf{v1} (\texttt{SatelliteNonGravToolbox}): Sentman shear $\times8$; air gas constant; spherical Earth (\SI{-21.4}{\kilo\metre}, density $\times1.40$); two solar constants; GPS as UT1, no ITRS$\to$GCRS; adjusted same-day $F_{10.7}$; silent density fallback; drag routine computes its own Sun; five optical implementations, torque without vis/IR split; single-epoch interfaces; hand-rolled tests; duplicated frame helpers. \textbf{v0.1.1} (\texttt{NonGravToolbox}): the physics specification is correct; the flaws are architectural: gravity and equations of motion inside a non-gravitational toolbox, a compatibility layer to a defective v1, a Sun ``selector'' that is a wrapper, a 4171-line handover mixing specification and diaries, a setup that does not add the tests, four instead of seven test categories, a 1367 oracle against a 1361 constant, no sensor layer, no quality control, error handling by convention only. Concept 3 keeps v0.1.1's dependency discipline, principles, help template and $\circ$/$\boxtimes$ rule, removes the scope creep, and adds the sensor and QC layers.

\section{Recommended defaults and test oracles}
\begin{table}[h]\centering\footnotesize
\begin{tabular}{@{}lp{6.6cm}p{5cm}@{}}\toprule
Quantity & Default & Authority\\\midrule
density & NRLMSISE-00 (\texttt{standard}), NRLMSIS 2.1 (\texttt{precise}) & Picone 2002; Emmert 2021\\
wind & co-rotation; HWM14 opt-in & Drob 2015\\
GSI & Sentman, fitted $\alpha_E=0.85$ in every preset & Sentman 1961; March 2021\\
$S_0$ & \SI{1361}{W/m^2}; 1367 override & IAU 2015 B3; SORCE/TIM\\
Lambert factor & $2/3$ & standard\\
shadow & conical; SOLAARS-CF in \texttt{precise} & Montenbruck \& Gill 2000; Robertson 2015\\
Earth radiation & Knocke $12\times24$; CERES $24\times48$ & Knocke 1988; Wielicki 1996\\
TRP & steady-state per panel; transient opt-in & Wöske 2019\\
magnetic field & IGRF-14 & IAGA 2024\\
ephemeris & SOFA analytic; DE440 in \texttt{precise} & \texttt{referenceSystem}\\
ellipsoid & WGS-84, passed explicitly & \\
GRACE-FO mass, geometry & \SI{601.214}{kg}; JPL D-56935 Tables 4--5 & Wen 2019\\\bottomrule
\end{tabular}
\end{table}
\begin{table}[h]\centering\footnotesize
\begin{tabular}{@{}p{9.6cm}p{6cm}@{}}\toprule
Oracle & Used for\\\midrule
CODATA 2018, IAU 2015 B3 & \texttt{const} (verified)\\
JPL D-56935 Tables 4--5 & \texttt{geom.graceFO} (verified: areas, optics, inertia, normals)\\
numpy Hampel (\texttt{tools/pyxcheck}) & \texttt{qc.hampel} (verified, bit-for-bit flags)\\
Brodowski NRLMSISE-00 DOCUMENTATION table (17 cases) & density backend\\
NRL/CCMC NRLMSIS 2.x and HWM14 test outputs; CNES DTM2020; SET JB2008 & extended density and wind\\
igrf14 test cases; Zenodo 10.5281/zenodo.14012303 & magnetic field\\
Montenbruck \& Gill \S3.4.2 example; Python disc integration & shadow function\\
Knocke 1988 Table 1; Orekit \texttt{KnockeRediffusedForceModel} & Earth radiation\\
Python kinetic-theory integration; v1 $\times8$ regression & Sentman coefficients\\
Pilinski 2013 Fig.~3 & SESAM\\
Siemes 2023 Table 3; Vielberg \& Kusche 2020 monthly means & radiation magnitudes per axis\\
Hu 2022 CoM calibration; Kornfeld 2019 ASD & electrostatic sensor\\
Lévèque 2022 CARIOQA parameters & atom-interferometer sensor\\
real ACT1B/GNV1B/SCA1B/THR1B/MAS1B 2020-01-01 GRACE-C & readers, end-to-end calibration (scale $\approx1.05$)\\\bottomrule
\end{tabular}
\end{table}
